\documentclass[11pt]{article}
\usepackage[T1]{fontenc}
\usepackage{textcomp}
\usepackage[margin=0.75in]{geometry}
\usepackage{amsmath,amssymb,amsthm}
\usepackage{array,booktabs}
\usepackage{hyperref}
\setlength{\parindent}{0pt}
\newtheorem{theorem}{Theorem}
\theoremstyle{definition}
\newtheorem{definition}{Definition}
\title{Flow Proof Artifact: Ring}
\author{Generated by \texttt{flow doc proof}}
\date{Flow Proof Book}
\begin{document}
\maketitle
Ring axioms extending commutative addition and associative multiplication.\\[0.5em]
\textbf{Source.} dummit-foote — *Abstract Algebra*, §7.1\\[0.5em]
\bigskip
\noindent{\large \textbf{Definition 1.} \textit{Ring addition commutes} \hfill Ring \cdot addition \cdot order does not matter}\par
\smallskip\noindent\textit{Coordinate: Ring \cdot addition \cdot order does not matter}\par
\smallskip\noindent\textit{Source: dummit-foote}\par
\medskip
\noindent\textbf{Goal.} Ring addition commutes\par
\noindent$\displaystyle \forall a \in Ring \forall b \in Ring\quad a + b = b + a$\par
\medskip
\noindent
\begin{tabular}{@{} >{\raggedright\arraybackslash}p{0.06\textwidth} >{\raggedright\arraybackslash}p{0.47\textwidth} >{\raggedleft\arraybackslash}p{0.41\textwidth} @{}}
\textbf{\#} & \textbf{Proof} & \textbf{Mathematics} \\
\midrule
\textbf{1.} & We stipulate order does not matter for addition on Ring — this is a definition, not a derived fact. & \\[0.45em]
\textbf{2.} & This follows directly from the definition: a plus b equals b plus a. Hence proven. & $\displaystyle a + b = b + a$ \\[0.45em]
\bottomrule
\end{tabular}
\label{thm:Ring-addition-order_does_not_matter}
\par\medskip\hrule
\bigskip
\noindent{\large \textbf{Definition 2.} \textit{Zero is the left additive identity} \hfill Ring \cdot addition \cdot zero is the left identity}\par
\smallskip\noindent\textit{Coordinate: Ring \cdot addition \cdot zero is the left identity}\par
\smallskip\noindent\textit{Source: dummit-foote}\par
\medskip
\noindent\textbf{Goal.} Zero is the left additive identity\par
\noindent$\displaystyle \forall r \in Ring\quad 0 + r = r$\par
\medskip
\noindent
\begin{tabular}{@{} >{\raggedright\arraybackslash}p{0.06\textwidth} >{\raggedright\arraybackslash}p{0.47\textwidth} >{\raggedleft\arraybackslash}p{0.41\textwidth} @{}}
\textbf{\#} & \textbf{Proof} & \textbf{Mathematics} \\
\midrule
\textbf{1.} & We stipulate zero is the left identity for addition on Ring — this is a definition, not a derived fact. & \\[0.45em]
\textbf{2.} & This follows directly from the definition: 0 plus r equals r. Hence proven. & $\displaystyle 0 + r = r$ \\[0.45em]
\bottomrule
\end{tabular}
\label{thm:Ring-addition-zero_is_the_left_identity}
\par\medskip\hrule
\bigskip
\noindent{\large \textbf{Definition 3.} \textit{One is the left multiplicative identity} \hfill Ring \cdot multiplication \cdot one is the left identity}\par
\smallskip\noindent\textit{Coordinate: Ring \cdot multiplication \cdot one is the left identity}\par
\smallskip\noindent\textit{Source: dummit-foote}\par
\medskip
\noindent\textbf{Goal.} One is the left multiplicative identity\par
\noindent$\displaystyle \forall r \in Ring\quad 1 \cdot r = r$\par
\medskip
\noindent
\begin{tabular}{@{} >{\raggedright\arraybackslash}p{0.06\textwidth} >{\raggedright\arraybackslash}p{0.47\textwidth} >{\raggedleft\arraybackslash}p{0.41\textwidth} @{}}
\textbf{\#} & \textbf{Proof} & \textbf{Mathematics} \\
\midrule
\textbf{1.} & We stipulate one is the left identity for multiplication on Ring — this is a definition, not a derived fact. & \\[0.45em]
\textbf{2.} & This follows directly from the definition: 1 times r equals r. Hence proven. & $\displaystyle 1 \cdot r = r$ \\[0.45em]
\bottomrule
\end{tabular}
\label{thm:Ring-multiplication-one_is_the_left_identity}
\par\medskip\hrule
\bigskip
\noindent{\large \textbf{Definition 4.} \textit{Multiplication distributes over addition on the left} \hfill Ring \cdot multiplication \cdot left distribution over addition holds}\par
\smallskip\noindent\textit{Coordinate: Ring \cdot multiplication \cdot left distribution over addition holds}\par
\smallskip\noindent\textit{Source: dummit-foote}\par
\medskip
\noindent\textbf{Goal.} Multiplication distributes over addition on the left\par
\noindent$\displaystyle \forall a \in Ring \forall b \in Ring \forall c \in Ring\quad a * (b + c) = a \cdot b + a \cdot c$\par
\medskip
\noindent
\begin{tabular}{@{} >{\raggedright\arraybackslash}p{0.06\textwidth} >{\raggedright\arraybackslash}p{0.47\textwidth} >{\raggedleft\arraybackslash}p{0.41\textwidth} @{}}
\textbf{\#} & \textbf{Proof} & \textbf{Mathematics} \\
\midrule
\textbf{1.} & We stipulate left distribution over addition holds for multiplication on Ring — this is a definition, not a derived fact. & \\[0.45em]
\textbf{2.} & This follows directly from the definition: a times (b plus c) equals a times b plus a times c. Hence proven. & $\displaystyle a * (b + c) = a \cdot b + a \cdot c$ \\[0.45em]
\bottomrule
\end{tabular}
\label{thm:Ring-multiplication-left_distribution_over_addition_holds}
\par\medskip\hrule
\end{document}
